\section{Introducción: aprender en entornos dinámicos}

Esta clase es la quinta de la serie de MIT 6.S191 (Introduction to Deep Learning), con el tema de \textbf{Deep Reinforcement Learning} (aprendizaje por refuerzo profundo). En las cuatro clases anteriores, el curso ya cubrió las arquitecturas centrales del aprendizaje profundo (redes completamente conectadas, RNN, CNN) y dos paradigmas de aprendizaje (Supervised Learning y Unsupervised Learning). Esta clase introduce el tercer paradigma de aprendizaje —\textbf{Reinforcement Learning} (aprendizaje por refuerzo)— y explora cómo combinar las técnicas de aprendizaje profundo con el aprendizaje por refuerzo para que un modelo pueda aprender una política óptima de comportamiento en un \textbf{entorno dinámico} mediante la interacción con dicho entorno.

\begin{figure}[H]
\centering
\includegraphics[width=0.85\textwidth]{slides-images/slide-02.jpg}
\caption{Aprender en entornos dinámicos: el aprendizaje por refuerzo permite que un agente aprenda a tomar decisiones a través de la interacción con el entorno\protect\footnotemark}
\end{figure}
\footnotetext{Diapositiva 2.}

La idea central del aprendizaje por refuerzo profundo difiere esencialmente de los paradigmas de aprendizaje anteriores: ya no depende de un conjunto de datos estático recolectado de antemano, sino que hace que el modelo (Agent) interactúe de manera continua dentro de un \textbf{entorno dinámico} (Environment) —tomando acciones (Action), observando los cambios del entorno (Observation) y recibiendo señales de recompensa (Reward)— para así aprender a tomar decisiones óptimas. Esta forma de aprendizaje tiene un amplio valor de aplicación en ámbitos como el control de robots, las estrategias de juego y la conducción autónoma.

\begin{knowledgebox}{¿Por qué se necesita el aprendizaje por refuerzo?}
En el mundo real, muchas tareas no pueden reducirse simplemente a un problema de aprendizaje supervisado del tipo «dada una entrada, predecir una salida». Por ejemplo, hacer que un robot aprenda a caminar o que una IA aprenda a jugar ajedrez requiere que el agente tome decisiones sucesivas a lo largo de una serie de pasos temporales, y la decisión actual afecta los estados y las recompensas futuros. Este es precisamente el escenario central de aplicación del aprendizaje por refuerzo.
\end{knowledgebox}

\begin{figure}[H]
\centering
\includegraphics[width=0.85\textwidth]{slides-images/slide-03.jpg}
\caption{Ámbitos de aplicación típicos del aprendizaje por refuerzo: robótica y estrategias de juego\protect\footnotemark}
\end{figure}
\footnotetext{Diapositiva 3.}

